\section{From Fixed Time Approximation to Hitting Time Mixing}\label{sec: fixed to hitting}

\subsection{Reduction to distributional approximation at a fixed time}
The role of this subsection is to upgrade the fixed-time approximation of \Cref{sec: Jain-Sawhney} into a conditional approximation given the exact set of untouched cards. This will provide the initial distribution for the marking argument started at $t^*$ defined in \eqref{eq: starting time of marking scheme}.

For a subset $T \subset [n]$, let $\mathcal{G}(T) = \mathcal{G}^t(T)$ denote the event that the random walk does not touch $T$ up to time $t$ and let $\mathcal{F}(T) = \mathcal{F}^t(T)$ denote the event that the random walk touches everything outside $T$ and does not touch $T$ up to time $t$. We use $\mathcal{F}(M)$ to denote $\mathcal{F}([M])$ for an integer $1\leq M \leq n$. Recall from \Cref{thm: X and nu} that $t =  \frac{n \log n}{k}  + t'$ where $|t'| \leq \frac{n}{2k}\log \log \log n$.


\begin{lemma}\label{lemma::Self-reduction lemma}
    Let $M \in \mathbb{N}$ and $T \subset [n]$ such that $M \leq |T| = o((\log n)^{3/2})$. Recall that $\gamma_{n,t} = \exp(-kt'/n).$ Then
    \begin{enumerate}
        \item $\mathbf{P}[\mathcal{G}(T)] = (1+(\log n)^{-1/2+o(1)})\mathbf{P}[\mathcal{G}([M])]\gamma_{n,t}^{|T| - M}n^{M-|T|}$,
        \item $\mathbf{P}[\mathcal{F}([M])] =(1+(\log n)^{-1/2+o(1)})\mathbf{P}[\mathcal{G}([M])]\exp(-\gamma_{n,t})$,
        \item $\sum_{T \supseteq [M], |T| \leq o((\log n)^{3/2})} \mathbf{P}[\mathcal{G}(T)] \leq (\log n)^{o(1)}\mathbf{P}[\mathcal{F}([M])]$.
    \end{enumerate}
\end{lemma}

\begin{proof}
For the first item,
\begin{equation}
\begin{aligned}
        \mathbf{P}[\mathcal{G}(T)] = \left(\frac{(n-|T|)_k}{(n)_k}\right)^{t} &= \exp\left(t\sum_{i=0}^{k-1}\log\left(1-\frac{|T|}{n-i}\right)\right)\\
        &= \exp\left(-kt\frac{|T|}{n}\right)\exp\left(O\left(\frac{tk^2|T|}{n^2}\right)+O\left(\frac{tk|T|^2}{n^2}\right)\right).
\end{aligned}
\end{equation}
Since $O\left(\frac{tk^2|T|}{n^2}\right) + O\left(\frac{tk|T|^2}{n^2}\right)=O\left(\frac{1}{(\log n)^{1/2}}\right)$, we have
\begin{equation}
\begin{aligned}
        \mathbf{P}[\mathcal{G}(T)]
        &= (1+(\log n)^{-1/2+o(1)})\exp\left(-kt\frac{|T|}{n}\right)\\
        &= (1+(\log n)^{-1/2+o(1)})\exp\left(-kt\frac{M}{n}\right)\exp\left(-kt\frac{|T|-M}{n}\right)\\
        &=  (1+(\log n)^{-1/2+o(1)})\mathbf{P}[\mathcal{G}(M)]\gamma_{n,t}^{|T|-M}n^{M-|T|}.
\end{aligned}
\end{equation}

For the second item, we use the principle of inclusion-exclusion to rewrite the indicator function of $\mathcal{F}(M)$ as a sum of indicator functions of $\mathcal{G}(T)$,
\begin{equation}
    \indic \{\mathcal{F}(M)\} = \sum_{T \supseteq [M]} (-1)^{|T| - M} \indic \{\mathcal{G}(T)\}.
\end{equation}
By the Bonferroni inequality, we have
\begin{equation}\label{eq: Bonferroni}
\sum_{T \supseteq [M], |T| \leq M + 2L +1} (-1)^{|T| - M} \indic \{\mathcal{G}(T)\} \leq \indic\{\mathcal{F}(M)\} \leq \sum_{T \supseteq [M], |T| \leq M + 2L} (-1)^{|T| - M} \indic \{\mathcal{G}(T)\}
\end{equation}
for any positive integer $L$. Take \(L=\lfloor\log n\rfloor\) and take expectations on both sides of \eqref{eq: Bonferroni}. Since \(\mathbf P[\mathcal G(T)]\) depends only on \(|T|\), the contribution of all sets \(T\supseteq[M]\) with \(|T|=M+j\) is
\(\binom{n-M}{j}\mathbf P[\mathcal G(T)]\), where \(T\) is any such set. Hence
$$
\sum_{j=0}^{2L+1}
(-1)^j
\binom{n-M}{j}
\mathbf P[\mathcal G(T)]
\leq
\mathbf P[\mathcal F(M)]
\leq
\sum_{j=0}^{2L}
(-1)^j
\binom{n-M}{j}
\mathbf P[\mathcal G(T)].
$$
Here, in the term indexed by \(j\), \(T\) denotes any set satisfying \(T\supseteq[M]\) and \(|T|=M+j\).
We then apply the first item
\[
\begin{aligned}
(1+(\log n)^{-1/2+o(1)})\frac{\mathbf{P}[\mathcal{F}(M)]}{\mathbf{P}[\mathcal{G}(M)]} &= \sum_{j=0}^{2L} \binom{n-M}{j}\frac{(-\gamma_{n,t})^j}{n^j} \pm  \binom{n-M}{2L+1}\frac{(-\gamma_{n,t})^{2L+1}}{n^{2L+1}}\\
&=(1+n^{-1+o(1)})\left(\sum_{j=0}^{2L} \frac{(-\gamma_{n,t})^j}{j!} \pm \frac{(-\gamma_{n,t})^{2L+1}}{(2L+1)!}  \right)\\
&= (1+n^{-1+o(1)})\exp(-\gamma_{n,t})
\end{aligned}
\]

where the last equality follows from the fact that $\gamma_{n,t} \leq \sqrt{\log \log n}$ and $L=\lfloor\log n\rfloor$.

For the last item, we use the first and second items
\[
\begin{aligned}
    \sum_{T \supseteq [M], T \leq o((\log n)^{3/2})} \mathbf{P}[\mathcal{G}(T)] &= \sum_{s=0}^{o((\log n)^{3/2})}\binom{n-M}{s}\mathbf{P}\left[\mathcal{G}([M+s])\right]\\
    & = (1+(\log n)^{-1/2+o(1)}) \mathbf{P}[\mathcal{F}(M)]e^{\gamma_{n,t}} \sum_{s=0}^{ o((\log n)^{3/2})} \binom{n-M}{s}\left( \frac{\gamma_{n,t}}{n}\right)^s\\
    & \leq (1+(\log n)^{-1/2+o(1)}) \mathbf{P}[\mathcal{F}(M)]e^{\gamma_{n,t}} \sum_{s=0}^{ o((\log n)^{3/2})} \binom{n}{s}\left( \frac{\gamma_{n,t}}{n}\right)^s\\
    &\leq (1+(\log n)^{-1/2+o(1)}) \mathbf{P}[\mathcal{F}(M)]e^{\gamma_{n,t}} \sum_{s=0}^{ o((\log n)^{3/2})}\left( \frac{\gamma_{n,t}^s}{s!}\right)\\
    &\leq (1+(\log n)^{-1/2+o(1)}) \mathbf{P}[\mathcal{F}(M)]e^{2\gamma_{n,t}}\\
    &\leq (\log n)^{o(1)}\mathbf{P}[\mathcal{F}(M)]
\end{aligned}
\]
where the last equality follows from the fact that $\gamma_{n,t} \leq \sqrt{\log \log n}$. 
\end{proof}

For the next lemma, suppose that $|t'| \leq \frac{n}{2k}\log \log \log n$. For a subset $T \subset [n]$, let 
\[
t_T := \frac{(n-|T|)\log(n-|T|)}{k}+ \frac{t'(n-|T|)}{n}.
\]
Then, we have 
\begin{equation}\label{eq: gamma_T}
\begin{aligned}
    \gamma_{n-|T|,t_T} = \exp(-kt'(n-|T|)/n(n-|T|)) = \exp(-kt'/n) =\gamma_{n,t}.
\end{aligned}
\end{equation}
\begin{defi}\label{def: nu^T_t}
    Let $\nu^T_t$ denote the distribution on $\mathfrak{S}_{[n]\setminus T}$  defined by the following sampling process. First, sample $M_t\in\{0,1,\ldots,n-|T|\}$ according to the distribution $\mathbf{P}[M_t=x]=\mathbf{P}(\Pois(\gamma_{n-|T|,t_T})=x)/\mathbf{P}(\Pois(\gamma_{n-|T|,t_T})\le n-|T|)$, then sample a uniformly random subset $S_t$ of size $M_t$ in $[n]\setminus T$. Finally,
\begin{enumerate}
\item When $k$ is odd, sample a uniformly random element of $\mathfrak{A}_{[n]\backslash S_t \cup T}$ and view it as an element of $\mathfrak{A}_n$ by fixing all of the elements in $S_t$ and $T$.
\item When $k$ is even, and $t$ is odd, sample a uniformly random element of $\mathfrak{A}_{[n]\backslash S_t\cup T}^c$ and view it as an element of $\mathfrak{A}_n^c$ by fixing all of the elements in $S_t$ and $T$.
\item When $k$ is even, and $t$ is even, sample a uniformly random element of $\mathfrak{A}_{[n]\backslash S_t\cup T}$ and view it as an element of $\mathfrak{A}_n$ by fixing all of the elements in $S_t$ and $T$.
\end{enumerate}
\end{defi}
Using natural inclusion from $\mathfrak{S}_{[n]\setminus T}$ to $\mathfrak{S}_n$, we view $\nu_t^T$ as a distribution on $\mathfrak{S}_n$. For the following \Cref{lem: relation of T and M} and \Cref{lem: conditioned}, we choose $ t = t_T$ and use the notation $\nu^T = \nu_{t_T}^T$. For the following lemma, we let $\Omega_{t_T}$ denote the parity class at the relevant actual time $t_T.$
\begin{lemma}\label{lem: relation of T and M}
    For any $\sigma \in \Omega_{t_T}$ such that $\mathrm{Fix}(\sigma) \supseteq [M]$ and $|\mathrm{Fix}(\sigma)| \leq {o((\log n)^{3/2})}$, we have
    \[
    \sum_{[M] \subseteq T \subseteq \mathrm{Fix}(\sigma)}(-1)^{|T| - M}\mathbf{P}[\mathcal{G}(T)]\nu^T(\sigma) = (2+(\log n)^{-1/2+o(1)})\frac{\mathbf{P}[\mathcal{F}(M)]}{(n-M)!}.
    \]
\end{lemma}

\begin{proof}
Notice that $\nu^T(\sigma) = 0$ if $T \not\subseteq \mathrm{Fix}(\sigma)$. Thus, we can only consider $\sigma \in \Omega_{t_T}$ such that $[M], T \subseteq \mathrm{Fix}(\sigma)$ and $|\mathrm{Fix}(\sigma)| \leq o((\log n)^{3/2})$, we see that
    \[
    \nu^T(\sigma) = \sum_{r=0}^{|\mathrm{Fix}(\sigma)| - |T|}\frac{\mathbf{P}(\mathrm{Pois}(\gamma_{n-|T|, t_T}) = r)}{\mathbf{P}(\mathrm{Pois}(\gamma_{n-|T|, t_T}) \leq n-|T|)}\frac{2\binom{|\mathrm{Fix}(\sigma)| -|T|}{r}}{\binom{n-|T|}{r}(n-|T|-r)!}.
    \]
Because $\gamma_{n-|T|, t_T} \leq \sqrt{\log \log n}$ and $|T| \leq o((\log n)^{3/2})$, we have that
\[
\mathbf{P}(\mathrm{Pois}(\gamma_{n-|T|, t_T}) > n-|T|) \leq \left( \frac{e\sqrt{\log \log n}}{n-|T|} \right)^{n-|T|} = O(n^{-1}).
\]
Also,
\[
\frac{\mathbf{P}(\mathrm{Pois}(\gamma_{n-|T|, t_T}) = r)}{n^{|T|-M}\binom{n-|T|}{r}(n-|T|-r)!} = \frac{e^{-\gamma_{n-|T|, t_T}}\gamma_{n-|T|, t_T}^r}{r!(n-|T|-r)!\binom{n-|T|}{r}n^{|T|-M}} = (1+n^{-1+o(1)})\frac{e^{-\gamma_{n-|T|, t_T}}\gamma_{n-|T|, t_T}^r}{(n-M)!}.
\]
Thus,
\[
\nu^T(\sigma) = (1+n^{-1+o(1)})\frac{e^{-\gamma_{n-|T|, t_T}}n^{|T|-M}}{(n-M)!} \sum_{r=0}^{|\mathrm{Fix}(\sigma)| - |T|} 2\gamma_{n-|T|, t_T}^r{\binom{|\mathrm{Fix}(\sigma)| -|T|}{r}}.
\]
Apply the first and second item of Lemma \ref{lemma::Self-reduction lemma} to the following expression, we have
\[
\begin{split}
    &(1+n^{-1+o(1)})(n-M)!\sum_{[M] \subseteq T \subseteq \mathrm{Fix}(\sigma)}(-1)^{|T| - M}\mathbf{P}[\mathcal{G}(T)]\nu^T(\sigma) \\
    &= (1+(\log n)^{-1/2+o(1)})\mathbf{P}[\mathcal{F}(M)]\sum_{[M] \subseteq T \subseteq \mathrm{Fix}(\sigma)}(-\gamma_{n-|T|, t_T})^{|T| - M}\sum_{r=0}^{|\mathrm{Fix}(\sigma)| - |T|} 2 \gamma_{n-|T|, t_T}^{r}\binom{|\mathrm{Fix}(\sigma)| -|T|}{r}\\
    & = (2+(\log n)^{-1/2+o(1)})\mathbf{P}[\mathcal{F}(M)]\sum_{[M] \subseteq T \subseteq \mathrm{Fix}(\sigma)}(-\gamma_{n-|T|, t_T})^{|T| - M}(1+\gamma_{n-|T|, t_T})^{|\mathrm{Fix}(\sigma)| - |T|}\\
    & = (2+(\log n)^{-1/2+o(1)})\mathbf{P}[\mathcal{F}(M)]\sum_{|T| - M \geq 0} \binom{|\mathrm{Fix}(\sigma)| - M}{|T|-M}
    (-\gamma_{n-|T|, t_T})^{|T|-M}(1+\gamma_{n-|T|, t_T})^{|\mathrm{Fix}(\sigma)| - |T|}\\
    &= (2+(\log n)^{-1/2+o(1)})\mathbf{P}[\mathcal{F}(M)].
\end{split}
\]
Here, we used binomial theorem in the second and fourth lines.
\end{proof}

\begin{lemma}\label{lem: conditioned}
Fix $t$ such that $|t'|\leq \frac{n}{k}(\frac{1}{2}\log\log\log n-2\log\log\log\log n)$. For $T \subset [n]$, let $\tilde{X}_t^{T}$ denote the random variable $X_t$ conditioned on the event $\mathcal{G}(T)$. Let $\nu_t^T$ denote the distribution described in \Cref{def: nu^T_t}. If $|T| = n^{o(1)}$, then
\[
d_{\mathrm{TV}}(\tilde{X}_t^T, \nu_t^T) \leq o\left((\log n)^{-5/2}\right).
\]
\end{lemma}

\begin{proof}
    Conditional on the event $\mathcal{G}(T)$, every sampled $k$-cycle is supported on $[n]\setminus T$, and the induced process on $[n]\setminus T$ is exactly the random $k$-cycle walk on $n-|T|$ points, run for $t$ steps.
    On the one hand, recall from \eqref{eq: gamma_T} that our times $t,t_T$ are designed to satisfy
    $$
        \gamma_{n-|T|,t_T}
        =
        \exp\left(
        -\frac{k}{n-|T|}\cdot \frac{n-|T|}{n}t'
        \right)
        =
        \exp\left(-\frac{kt'}{n}\right)
        =
        \gamma_{n,t}.
    $$

    On the other hand, when we apply \Cref{thm: X and nu} to the conditioned walk on
    $[n]\setminus T$ at the actual time $t$, the corresponding time shift is
    $$
        t-\frac{(n-|T|)\log(n-|T|)}{k}.
    $$
    Since
    $$
        t-\frac{(n-|T|)\log(n-|T|)}{k}=t'+O\left(\frac{|T|\log n}{k}+1\right),
    $$
    the corresponding Poisson parameter differs from
    $$
        \exp\left(-\frac{kt'}{n-|T|}\right)
    $$
    by a multiplicative factor $1+n^{-1+o(1)}$. It remains to compare
    $\exp(-kt'/n)$ and $\exp(-kt'/(n-|T|))$. We have
    \begin{align}
        \left|
        \exp\left(-\frac{kt'}{n}\right)
        -
        \exp\left(-\frac{kt'}{n-|T|}\right)
        \right|
        &\leq
        \exp\left(\frac{k|t'|}{n}\right)
        \left|
        1-\exp\left(-\frac{kt'|T|}{n(n-|T|)}\right)
        \right| \notag\\
        &\leq
        2\exp\left(\frac{k|t'|}{n}\right)
        \frac{k|t'||T|}{n(n-|T|)} \notag\\
        &= n^{-1+o(1)}.
    \end{align}
    Therefore, by the standard coupling of Poisson random variables with nearby means, the total variation distance between the two corresponding auxiliary measures is $n^{-1+o(1)}$.

    It remains to check that the actual time $t$ lies in the window where \Cref{thm: X and nu}
    applies to the random $k$-cycle walk on $n-|T|$ points. We need
    \begin{equation}
        \left|
        t-\frac{(n-|T|)\log(n-|T|)}{k}
        \right|
        \leq
        \frac{n-|T|}{2k}\log\log\log (n-|T|).
    \end{equation}
    Since $t=\frac{n\log n}{k}+t'$, it suffices to compare $n\log n$ with $(n-|T|)\log(n-|T|)$. We have
    \begin{equation}
        \left|
        \frac{n\log n-(n-|T|)\log(n-|T|)}{k}
        \right|
        \leq
        O\left(\frac{|T|\log n}{k}\right)
        =
        o\left(\frac{n}{k}\log\log\log\log n\right),
    \end{equation}
    while
    \begin{equation}
        \frac{n-|T|}{2k}\log\log\log(n-|T|)
        =
        \frac{n}{2k}\log\log\log n
        +
        o\left(\frac{n}{k}\log\log\log\log n\right).
    \end{equation}
    Hence the assumed bound on $t'$ implies the desired window condition for $n-|T|$. Applying \Cref{thm: X and nu} to the conditioned walk on $[n]\setminus T$, and then using
    the above comparison of the Poisson parameters, gives
    $$
        d_{\mathrm{TV}}\left(\widetilde X_t^T,\nu_t^T\right)
        \leq
        o\left((\log n)^{-5/2}\right)+n^{-1+o(1)}
        =
        o\left((\log n)^{-5/2}\right).
    $$
    This completes the proof.
\end{proof}

For each integer \(t\), define the parity class
$$
\Omega_t:=
\begin{cases}
\mathfrak A_n, & k \text{ is odd}\\
\mathfrak A_n, & k \text{ is even and } t \text{ is even}\\
\mathfrak A_n^c, & k \text{ is even and } t \text{ is odd}
\end{cases}.
$$
For a subset \(M\subseteq [n]\), set
$$
    \Omega_t(M):=\{\sigma\in \Omega_t:\sigma(i)=i\text{ for every }i\in M\}.
$$
Let \(\mu_t^M\) denote the uniform measure on \(\Omega_t(M)\). Equivalently, \(\mu_t^M\) is obtained by sampling uniformly from the appropriate parity class on \([n]\setminus M\), and then extending the permutation to \([n]\) by fixing every point of \(M\).

In the following proposition, we show that when the time is close to $n\log n/k$, the random walk is nearly uniform after conditioning on the event that the set of untouched cards is exactly $M$.

\begin{prop}\label{prop:key}
    Fix $t$ and $M\subseteq [n]$ such that
    $$
        |t'|\leq \frac{n}{k}
        \left(
        \frac{1}{2}\log\log\log n
        -2\log\log\log\log n
        \right)
        \qquad\text{and}\qquad
        |M|\leq \log n.
    $$
    Let $X_t^M$ denote the random variable $X_t$ conditioned on the event
    $\mathcal F(M)$. Then
    $$
        d_{\mathrm{TV}}(X_t^M,\mu_t^M)
        \leq
        (\log n)^{-1/2+o(1)}.
    $$
\end{prop}

\begin{proof}
    By the definition of total variation distance,
    \begin{align*}
        d_{\mathrm{TV}}(X_t^M,\mu_t^M)
        &=
        \frac{1}{2}
        \sum_{\sigma\in \Omega_t(M)}
        \left|
        \mathbf P(X_t^M=\sigma)-\mu_t^M(\sigma)
        \right| \\
        &=
        \frac{1}{2}
        \sum_{\sigma\in \Omega_t(M)}
        \left|
        \frac{\mathbf P(\{X_t=\sigma\}\cap \mathcal F(M))}
        {\mathbf P(\mathcal F(M))}
        -
        \frac{2}{(n-|M|)!}
        \right|.
    \end{align*}
    Therefore, it suffices to prove
    $$
        \sum_{\sigma\in \Omega_t(M)}
        \left|
        \mathbf P(\{X_t=\sigma\}\cap \mathcal F(M))
        -
        \frac{2\mathbf P(\mathcal F(M))}{(n-|M|)!}
        \right|
        \leq
        (\log n)^{-1/2+o(1)}\mathbf P(\mathcal F(M)).
    $$

    We split $\Omega_t(M)$ according to the number of fixed points. Let
    $$
        S:=
        \Omega_t(M)\cap
        \{\sigma\in \Omega_t:|\operatorname{Fix}(\sigma)|\leq 3\log n\},
        \qquad
        L:=\Omega_t(M)\setminus S.
    $$
    We bound the contributions from $L$ and $S$ separately.

    \medskip
    \noindent
    \textbf{Contribution from $L$.}
    By \Cref{lem: conditioned}, we have
    $$
        \sum_{\sigma\in \Omega_t(M)}
        \left|
        \mathbf P(\widetilde X_t^M=\sigma)-\nu^M(\sigma)
        \right|
        =
        o\left((\log n)^{-5/2}\right).
    $$
    Recall that
    $$
        t_M=\frac{(n-|M|)\log(n-|M|)}{k}
        +\frac{n-|M|}{n}t',
    $$
    and that $\nu^M$ is obtained by first sampling a Poisson number of additional
    fixed points in $[n]\setminus M$, with parameter $\gamma_{n-|M|,t_M}$, and then
    sampling uniformly from the appropriate parity class on the remaining points.
    Hence
    \begin{align*}
        \sum_{\sigma\in L}\mathbf P(\widetilde X_t^M=\sigma)
        &\leq
        \sum_{\sigma\in L}
        \left|
        \mathbf P(\widetilde X_t^M=\sigma)-\nu^M(\sigma)
        \right|
        +\nu^M(L) \\
        &\leq
        o\left((\log n)^{-5/2}\right)
        +
        \mathbf P\left(\operatorname{Pois}(\gamma_{n-|M|,t_M})\geq \log n\right) +
        \sup_{\substack{Q\subseteq [n]\setminus M\\ |Q|\leq \log n}}
        \mu_t^{M\cup Q}(L).\\
    \end{align*}
    Since
    $$
        \gamma_{n-|M|,t_M}
        =
        \gamma_{n,t}
        \leq
        \exp\left(
        \frac{1}{2}\log\log\log n
        -2\log\log\log\log n
        \right)
        \leq
        (\log\log n)^{1/2+o(1)},
    $$
    Chernoff's bound gives
    $$
        \mathbf P\left(\operatorname{Pois}(\gamma_{n-|M|,t_M})\geq \log n\right)
        \leq
        \left(
        \frac{e\sqrt{\log\log n}}{\log n}
        \right)^{\log n}
        \leq
        (\log n)^{-1/2+o(1)}.
    $$
    Moreover, uniformly over all $Q\subseteq [n]\setminus M$ with $|Q|\leq \log n$, a uniform permutation on $[n]\setminus (M\cup Q)$ has at least $\log n$ fixed
    points with probability at most $
    \frac{2}{(\log n)!}\leq(\log n)^{-1/2+o(1)}$. Therefore,
    $$
        \sum_{\sigma\in L}\mathbf P(\widetilde X_t^M=\sigma)
        \leq
        (\log n)^{-1/2+o(1)}.
    $$

    Since $\mathcal F(M)\subseteq \mathcal G(M)$, we have
    \begin{align*}
        \sum_{\sigma\in L}
        \mathbf P(\{X_t=\sigma\}\cap \mathcal F(M))
        &\leq
        \sum_{\sigma\in L}
        \mathbf P(\{X_t=\sigma\}\cap \mathcal G(M)) \\
        &=
        \mathbf P(\mathcal G(M))
        \sum_{\sigma\in L}
        \mathbf P(\widetilde X_t^M=\sigma) \\
        &\leq
        (\log n)^{-1/2+o(1)}\mathbf P(\mathcal G(M)).
    \end{align*}
    By \Cref{lemma::Self-reduction lemma},
    $$
        \mathbf P(\mathcal G(M))
        \leq
        (\log n)^{o(1)}\mathbf P(\mathcal F(M)).
    $$
    Hence
    $$
        \sum_{\sigma\in L}
        \mathbf P(\{X_t=\sigma\}\cap \mathcal F(M))
        \leq
        (\log n)^{-1/2+o(1)}\mathbf P(\mathcal F(M)).
    $$
    On the other hand, since $\mu_t^M$ is uniform on $\Omega_t(M)$,
    $$
        \sum_{\sigma\in L}
        \frac{2\mathbf P(\mathcal F(M))}{(n-|M|)!}
        =
        \mathbf P(\mathcal F(M))\mu_t^M(L)
        \leq
        (\log n)^{-1/2+o(1)}\mathbf P(\mathcal F(M)).
    $$
    Combining the preceding two estimates gives
    $$
        \sum_{\sigma\in L}
        \left|
        \mathbf P(\{X_t=\sigma\}\cap \mathcal F(M))
        -
        \frac{2\mathbf P(\mathcal F(M))}{(n-|M|)!}
        \right|
        \leq
        (\log n)^{-1/2+o(1)}\mathbf P(\mathcal F(M)).
    $$

    \medskip
    \noindent
    \textbf{Contribution from $S$.}
    For $\sigma\in S$, the inclusion-exclusion identity gives
    $$
        \mathbf 1_{\mathcal F(M)}
        =
        \sum_{T\supseteq M}
        (-1)^{|T|-|M|}\mathbf 1_{\mathcal G(T)}.
    $$
    Since $\nu^T(\sigma)=0$ unless $T\subseteq \operatorname{Fix}(\sigma)$, we may restrict
    the sum to $M\subseteq T\subseteq \operatorname{Fix}(\sigma)$. By \Cref{lem: relation of T and M},
    for every $\sigma\in S$ we have
    $$
        \sum_{M\subseteq T\subseteq \operatorname{Fix}(\sigma)}
        (-1)^{|T|-|M|}
        \mathbf P(\mathcal G(T))\nu^T(\sigma)
        =
        \left(2+(\log n)^{-1/2+o(1)}\right)
        \frac{\mathbf P(\mathcal F(M))}{(n-|M|)!}.
    $$
    Therefore,
    \begin{align*}
        &\sum_{\sigma\in S}
        \left|
        \mathbf P(\{X_t=\sigma\}\cap \mathcal F(M))
        -
        \frac{2\mathbf P(\mathcal F(M))}{(n-|M|)!}
        \right| \\
        &\leq
        (\log n)^{-1/2+o(1)}\mathbf P(\mathcal F(M))+
        \sum_{\sigma\in S}
        \sum_{M\subseteq T\subseteq \operatorname{Fix}(\sigma)}
        \left|
        \mathbf P(\{X_t=\sigma\}\cap \mathcal G(T))
        -
        \mathbf P(\mathcal G(T))\nu^T(\sigma)
        \right| \\
        &\leq
        (\log n)^{-1/2+o(1)}\mathbf P(\mathcal F(M))\quad+
        \sum_{\substack{T\supseteq M\\ |T|\leq 3\log n}}
        \mathbf P(\mathcal G(T))
        \sum_{\sigma\in S}
        \left|
        \mathbf P(X_t=\sigma\mid \mathcal G(T))
        -
        \nu^T(\sigma)
        \right|. \\
    \end{align*}
    By \Cref{lem: conditioned}, for every $T$ appearing in the last sum,
    $$
        \sum_{\sigma\in S}
        \left|
        \mathbf P(X_t=\sigma\mid \mathcal G(T))
        -
        \nu^T(\sigma)
        \right|
        \leq
        2d_{\mathrm{TV}}(\widetilde X_t^T,\nu^T)
        =
        o\left((\log n)^{-5/2}\right).
    $$
    Hence, using the last item of \Cref{lemma::Self-reduction lemma},
    \begin{align*}
        &\sum_{\sigma\in S}
        \left|
        \mathbf P(\{X_t=\sigma\}\cap \mathcal F(M))
        -
        \frac{2\mathbf P(\mathcal F(M))}{(n-|M|)!}
        \right| \\
        &\leq
        (\log n)^{-1/2+o(1)}\mathbf P(\mathcal F(M))
        +
        o\left((\log n)^{-5/2}\right)
        \sum_{\substack{T\supseteq M\\ |T|\leq 3\log n}}
        \mathbf P(\mathcal G(T)) \\
        &\leq
        (\log n)^{-1/2+o(1)}\mathbf P(\mathcal F(M)).
    \end{align*}

    Combining the contributions from $S$ and $L$ proves the desired estimate, and hence
    completes the proof.
\end{proof}

\subsection{The marking schemes for odd and even $k$-cycles}\label{subsec: marking schemes}

In this subsection, we introduce the marking schemes used in the proof of \Cref{thm: hitting time mixing_main thm}. The role of these schemes is to turn the conditional uniformity statement from \Cref{prop:key} into a stopped-time statement. We start the marking procedure at the deterministic time
$t^*$ defined in \eqref{eq: starting time of marking scheme}. At this time, all cards which have already been touched are declared marked, and the remaining cards are marked one by one during the subsequent evolution.

The construction has to serve two purposes at once. On the one hand, it should remain close to the genuine touching process of the random $k$-cycle walk, so that the terminal marking time is close to the hitting time $\tau$. On the other
hand, it should be artificial enough to preserve an exact conditional uniformity property after conditioning on the current marked set. We therefore use two related marking schemes. Marking scheme B uses the original transition kernel $P_n$ and is coupled directly to the original walk. Marking scheme A uses a
modified time-inhomogeneous transition kernel $Q_{n,t}$, designed so that an auxiliary process remains conditionally uniform on the appropriate parity class.

More precisely, the proof will compare four processes
$$
\rho_t\longleftrightarrow Y_t\longleftrightarrow Z_t\longleftrightarrow X_t .
$$
Here $X_t$ is the original random $k$-cycle walk. The process $Z_t$ uses marking
scheme B and is therefore closest to $X_t$. The process $Y_t$ uses marking scheme
A and is coupled to $Z_t$ by comparing the kernels $Q_{n,t}$ and $P_n$.
Finally, $\rho_t$ uses the same $Q_{n,t}$-moves and the same marked sets as
$Y_t$, but its action on the permutation is modified in order to maintain exact
conditional uniformity. Thus $\rho_t$ is the process which supplies the uniform
distribution at the terminal marking time, while $Y_t$ and $Z_t$ bridge it back
to the original chain.

For the rest of this subsection, let $S_{t^*}$ denote the set of elements in
$[n]$ which have not been touched by time $t^*$. For $t\ge t^*$, let $S_t$
denote the set of unmarked elements and let
$$
M_t=[n]\setminus S_t
$$
be the marked set. The odd and even cases have the same basic marking structure:
a useful update is one which contains exactly one element of $S_t$, in which
case that element is marked. The only additional issue is parity. When $k$ is
odd, each $k$-cycle is an even permutation, so the walk stays in a fixed parity
class. When $k$ is even, each $k$-cycle is odd, so the auxiliary chain must also
change parity at every step. This is the reason for the parity-correcting
transpositions supported on the marked set in the definition of $Q_{n,t}$ below.

\begin{defi}[The modified transition kernel $Q_{n,t}$]
For $t\ge t^*$, set
$$
C_t=(|M_t|-1)_{k-2},
\qquad
q_t=\frac{|M_t|}{|M_t|+1}
\left(\frac{n!}{k(n-k)!}\right)^{-1}.
$$
If $k$ is odd, then the kernel $Q_{n,t}$ is defined by
\[
Q_{n,t}=\begin{cases}
(i\cdots i) & \text{with probability } C_tq_t \ \text{for } i\in S_t,\\[6pt]
(i_1\cdots i_k) & \text{with probability } q_t\ \text{for } i_1\in S_t,\ i_2,\ldots,i_k\in M_t,\\[6pt]
(i_1\cdots i_k) & \text{with probability } \left(\dfrac{n!}{k(n-k)!}\right)^{-1}\ \text{otherwise}.
\end{cases}
\]
If $k$ is even, then the kernel $Q_{n,t}$ is defined by
\[
Q_{n,t}=\begin{cases}
(ab)(i\cdots i) & \text{with probability } \dfrac{C_tq_t}{\binom{|M_t|}{2}}
\ \text{for } i\in S_t,\ \{a,b\}\subset M_t,\\[8pt]
(i_1\cdots i_k) & \text{with probability } q_t\ \text{for } i_1\in S_t,\ i_2,\ldots,i_k\in M_t,\\[6pt]
(i_1\cdots i_k) & \text{with probability } \left(\dfrac{n!}{k(n-k)!}\right)^{-1}\ \text{otherwise}.
\end{cases}
\]

Here $(i\cdots i)$ is a labelled dummy move. In the odd case it acts as the identity. In the even case $(ab)(i\cdots i)$ acts on the permutation as the transposition $(ab)$ and carries the label $i$, which indicates the element to be marked. Throughout the marking argument we restrict to the event $|S_{t^*}|\leq \log n$, whose complement will be handled separately. Under both marking schemes the number of unmarked elements never increases, so for all sufficiently large $n$ we have $|M_t|\geq n-\log n$ for every $t\geq t^*$. Hence all quantities appearing in the definition of $Q_{n,t}$ are well-defined.


Let us check that $Q_{n,t}$ is indeed a probability kernel. The total mass assigned by the original uniform $k$-cycle kernel to cycles with exactly one element in $S_t$ is
$$
|S_t|\cdot|M_t|\cdot(|M_t|-1)_{k-2}
\left(\frac{n!}{k(n-k)!}\right)^{-1}=|S_t|\cdot|M_t|\cdot C_t\cdot
\left(\frac{n!}{k(n-k)!}\right)^{-1}.
$$
Under $Q_{n,t}$, this mass is replaced by two contributions. The ordinary one-unmarked $k$-cycles have total mass $|S_t|\cdot|M_t|\cdot C_t q_t$,
while the dummy moves have total mass
$|S_t|\cdot C_t q_t$ in the odd case, and
$$
|S_t|\cdot \binom{|M_t|}{2}\cdot\frac{C_tq_t}{\binom{|M_t|}{2}}=|S_t|\cdot C_t q_t
$$
in the even case. Hence the total replacement mass is
$$
|S_t|\cdot (|M_t|+1)\cdot C_tq_t=|S_t|\cdot|M_t|\cdot C_t
\left(\frac{n!}{k(n-k)!}\right)^{-1},
$$
by the definition of $q_t$. All other $k$-cycles retain their original mass. Therefore the total mass of $Q_{n,t}$ is indeed one.

\end{defi}

\begin{defi}[Marking scheme A]
At time $t^*$, mark all elements in $[n]\setminus S_{t^*}$. At time $t+1$,
sample a move $\sigma$ from the kernel $Q_{n,t}$.

The marked set is updated as follows.
\begin{itemize}
\item If the sampled move is a dummy move labelled by $i\in S_t$, then mark
$i$.
\item If $\sigma=(i_1\cdots i_k)$ with $i_1\in S_t$ and
$i_2,\ldots,i_k\in M_t$, then mark $i_1$. Thus, $M_{t+1} = M_t \cup \{i_1\}$ and $S_{t+1} = S_t \setminus \{i_1\}$.
\item In all remaining cases, namely when the ordinary $k$-cycle does not
mark exactly one new element, pull back the marked set under $\sigma$:
$$
M_{t+1}=\{i\in[n]:\sigma(i)\in M_t\}=\sigma^{-1}(M_t),
\qquad
S_{t+1}=\{i\in[n]:\sigma(i)\in S_t\}=\sigma^{-1}(S_t).
$$
\end{itemize}
\end{defi}

\begin{defi}[Marking scheme B]
At time $t^*$, mark all elements in $[n]\setminus S_{t^*}$. At time $t+1$,
sample a genuine $k$-cycle
$$
\sigma=(i_1\cdots i_k)
$$
from the original transition kernel $P_n$.

The marked set is updated by the same rule as in marking scheme A, except that
there are no dummy moves. Thus, if exactly one point of the cycle is in $S_t$,
say $i_1\in S_t$ and $i_2,\ldots,i_k\in M_t$, then mark $i_1$ and let $M_{t+1} = M_t \cup \{i_1\}$ and $S_{t+1} = S_t \setminus \{i_1\}$. Otherwise, pull
back the marked set:
$$
M_{t+1}=\sigma^{-1}(M_t),
\qquad
S_{t+1}=\sigma^{-1}(S_t).
$$
\end{defi}

We now define the coupled processes which will be compared in the proof. Let
$$
\rho_{t^*}=Y_{t^*}=Z_{t^*}
$$
be a common random variable sampled uniformly from
$\Omega_{t^*}(S_{t^*})$.

\begin{defi}[The coupled processes]
We define four Markov processes
$$
(X_t)_{t\ge t^*},\qquad
(Y_t)_{t\ge t^*},\qquad
(Z_t)_{t\ge t^*},\qquad
(\rho_t)_{t\ge t^*}
$$
on a common probability space as follows.
\begin{itemize}
\item The process $(X_t)_{t\ge t^*}$ is the original random $k$-cycle walk,
evolving according to the transition kernel $P_n$. Let $\tau$ denote
the first time at which all elements have been touched.

\item The process $(Z_t)_{t\ge t^*}$ is initialized by $Z_{t^*}$ and evolves
according to the original transition kernel $P_n$. Its marked set is
updated according to marking scheme B. Let $\tau_1$ denote the first time at
which all elements have been marked.

\item The process $(Y_t)_{t\ge t^*}$ is initialized by $Y_{t^*}$ and evolves
according to the modified transition kernel $Q_{n,t}$. Its marked set is
updated according to marking scheme A. Let $\tau_2$ denote the first time at
which all elements have been marked.

\item The process $(\rho_t)_{t\ge t^*}$ is initialized by $\rho_{t^*}$ and
uses the same $Q_{n,t}$-moves and the same marked sets as $Y_t$. Its action
on the current state is modified as follows.

At time $t+1$, sample $\sigma\sim Q_{n,t}$.

\begin{itemize}
    \item If the sampled move marks exactly one new element $i$, either through
    a dummy move or through a one-unmarked $k$-cycle, set
    $$
      M_{t+1}:=M_t \cup \{i\},
        \qquad
        S_{t+1}:=S_t \setminus \{i\}, \qquad  \rho_{t+1}:=\sigma\rho_t.
    $$

    \item If no new element is marked and
    $\operatorname{supp}(\sigma)\subset M_t$, set
    $$
       M_{t+1}:=M_t,
        \qquad
        S_{t+1}:=S_t, \qquad  \rho_{t+1}:=\sigma\rho_t.
    $$

    \item If no new element is marked and
    $\operatorname{supp}(\sigma)\not\subset M_t$, first update
    $$
        M_{t+1}:=\sigma^{-1}(M_t),
        \qquad
        S_{t+1}:=\sigma^{-1}(S_t).
    $$
    If $k$ is odd, set
    $$
        \rho_{t+1}:=\sigma^{-1}\rho_t\sigma.
    $$
    If $k$ is even, then, conditionally on $M_{t+1}$, choose an unordered pair
    $\{U_{t+1},V_{t+1}\}\subset M_{t+1}$ uniformly and independently of
    $\rho_t$, and set
    $$
        \rho_{t+1}:=(U_{t+1}\,V_{t+1})\sigma^{-1}\rho_t\sigma.
    $$
\end{itemize}
Let $\tau_3$ denote the first time at which all elements have been marked
for the process $\rho_t$.

\end{itemize}
\end{defi}


\subsection{Proof of \Cref{thm: hitting time mixing_main thm}}
Finally, we are ready for the proof of \Cref{thm: hitting time mixing_main thm}. First, we show that, conditioned on the marked set \(M_t\), the random permutation \(\rho_t\) is uniform on \(\Omega_t(S_t)\). In particular, at the marking time \(\tau_3\), all elements are marked, and hence \(\rho_{\tau_3}\sim U_{\mathfrak A_n}\) when $k$ is odd, and   \(\rho_{\tau_3}\sim \mathbf{P}(\tau_3 ~\text{even})\,U_{\mathfrak A_n}+\mathbf{P}(\tau_3 ~\text{odd})\,U_{\mathfrak A_n^c}\) when $k$ is even.

Second, we compare the auxiliary chains with the original chain. More precisely, we show that
\[
    d_{\mathrm{TV}}(\rho_{\tau_3},Y_{\tau_2}),\qquad
    d_{\mathrm{TV}}(Y_{\tau_2},Z_{\tau_1}),\qquad
    d_{\mathrm{TV}}(Z_{\tau_1},X_{\tau})
\]
are all small. Combining these estimates with the uniformity of \(\rho_{\tau_3}\), together with the fact that in the even-\(k\) case the hitting time has asymptotically balanced parity, will yield the desired hitting-time mixing statement. The start coupling time $t^*$ will again be the same as in \eqref{eq: starting time of marking scheme}. The following lemma is the key reason for introducing the modified process $\rho_t$: it shows that the artificial dynamics preserve exact conditional uniformity, so that once all elements have been marked, the terminal distribution is the desired uniform measure on the appropriate parity class.

\begin{lemma}\label{lem:uniform}
    For $t \geq t^*$, conditioned on the marked set $M_t$, $\rho_t$ is uniform on $\Omega_{t}(S_t)$. In particular, $\rho_{\tau_3}|\{\tau_3= t\}$ is uniform on $\Omega_t$. Consequently, if $k$ is odd, then
    \begin{align*}
        \rho_{\tau_3}\sim U_{\mathfrak{A}_n}.
    \end{align*}
    If $k$ is even, then
\begin{align*}
    \rho_{\tau_3}\sim \mathbf{P}(\tau_3\text{ is even}) \, U_{\mathfrak{A}_n}+\mathbf{P}(\tau_3\text{ is odd}) \, U_{\mathfrak{A}_n^c}.
\end{align*}
\end{lemma}
\begin{proof}
    At time $t^*$, the claim follows from the construction of $\rho_{t^*}$. Suppose that the claim holds at time $t\geq t^*$. We prove it at time $t+1$. 


First consider the case in which no new element is marked. There are two subcases. If the sampled $k$-cycle satisfies $\operatorname{supp}(\sigma)\subset M_t$, then $M_{t+1} = M_t, S_{t+1}= S_t,\rho_{t+1}= \sigma\rho_t$. Left multiplication by $\sigma$ gives a bijection $\Omega_t(S_t)\longrightarrow \Omega_{t+1}(S_{t+1})$. Hence $\rho_{t+1}$ is uniform on $\Omega_{t+1}(S_{t+1})$. On the other hand, if the sampled $k$-cycle satisfies $\operatorname{supp}(\sigma)\not\subset M_t$, then $M_{t+1} = \sigma^{-1}(M_t), S_{t+1}= \sigma^{-1}(S_t)$. If $k$ is odd, then $\rho_{t+1}= \sigma^{-1}\rho_t\sigma$. For every sampled move $\sigma$, conjugation by $\sigma$ gives a bijection $\Omega_{t}(S_t)\rightarrow \Omega_{t+1}(S_{t+1})$. Therefore, $\rho_{t+1}$ is uniform on $\Omega_{t+1}(S_{t+1})$. If $k$ is even, then the random permutation
$
\sigma^{-1}\rho_t\sigma
$
is uniform on
$
\sigma^{-1}\Omega_t(S_t)\sigma=\Omega_t(S_{t+1}).
$ Conditionally on $M_{t+1}$, the unordered pair
$
\{U_{t+1},V_{t+1}\}
$
is chosen uniformly from the set of unordered pairs contained in $M_{t+1}$, independently of
$\rho_t$, and
\[
\rho_{t+1}=(U_{t+1}V_{t+1})\sigma^{-1}\rho_t\sigma.
\]
For every fixed unordered pair $\{u,v\}\subset M_{t+1}$, left multiplication by the transposition
$(uv)$ gives a bijection
$
\Omega_t(S_{t+1})\longrightarrow \Omega_{t+1}(S_{t+1}),
$
because $(uv)$ is supported on $M_{t+1}$ and changes parity. Averaging over the uniformly
chosen pair therefore preserves the uniform law on $\Omega_{t+1}(S_{t+1})$.


It remains to consider the case in which exactly one new element is marked. Then there exists
$i\in S_t$ such that
\[
M_{t+1}=M_t\cup\{i\},\qquad S_{t+1}=S_t\setminus\{i\}.
\]
Fix $\pi\in \Omega_{t+1}(S_{t+1})$. Since $\pi$ fixes every element of
$S_t\setminus\{i\}$, bijectivity implies $\pi(i)\in M_t\cup\{i\}$. 
We show
that the total transition weight of all moves taking a state in
$\Omega_t(S_t)$ to $\pi$ is independent of $\pi$.

Let $\alpha$ denote the sampled move acting on $\rho_t$, and write
$m=|M_t|$. If $k$ is odd, the moves that mark $i$ are the identity dummy move,
with weight $C_tq_t$, and the one-unmarked cycles
$(i\,a_2\,\cdots\,a_k)$ with $a_2,\ldots,a_k\in M_t$ distinct, each with
weight $q_t$. If $\pi(i)=i$, only the identity dummy move contributes, so the
total contribution is $C_tq_t=(m-1)_{k-2}q_t$. If $\pi(i)\in M_t$, then a
contributing genuine cycle must satisfy $\alpha(i)=\pi(i)$; hence
\[
    \alpha=(i\,\pi(i)\,a_3\,\cdots\,a_k),
\]
where $a_3,\ldots,a_k$ are distinct elements of
$M_t\setminus\{\pi(i)\}$. There are $(m-1)_{k-2}=C_t$ such cycles, so the total
contribution is again $C_tq_t$.

If $k$ is even, the moves that mark $i$ are the dummy transpositions
$(ab)(i\cdots i)$, where $\{a,b\}\subset M_t$, each with weight
$C_tq_t/\binom{m}{2}$, and the same one-unmarked genuine cycles as above, each
with weight $q_t$. If $\pi(i)=i$, only dummy transpositions can contribute, and
their total contribution is
\[
    \binom{m}{2}\cdot \frac{C_tq_t}{\binom{m}{2}}=C_tq_t.
\]
If $\pi(i)\in M_t$, no dummy transposition can contribute because every dummy
transposition fixes $i$. The same counting of genuine cycles as in the odd case
therefore gives total contribution $C_tq_t$.

Thus, in both parity cases, the total preimage weight is independent of the
target permutation $\pi\in\Omega_{t+1}(S_{t+1})$. Consequently, conditionally on
marking the element $i$, the random permutation $\rho_{t+1}$ is uniform on
$\Omega_{t+1}(S_{t+1})$.
\end{proof} 

%We next compare the artificial process $\rho_t$ with the process $Y_t$, which uses the same marked sets but evolves by ordinary left multiplication.

The next lemma compares marking scheme A with marking scheme B by showing that the corresponding terminal times and stopped permutations agree with high probability under a natural coupling.

\begin{lemma}\label{lem: Z and Y}
    With notations defined in \Cref{subsec: marking schemes}, we have $\mathbf P(\tau_1\neq \tau_2)\le n^{-1+o(1)}$
    and
    $$d_{\mathrm{TV}}(\mathcal{L}(Z_{\tau_1}),\mathcal{L}(Y_{\tau_2})) \leq n^{-1+o(1)}.$$
\end{lemma}

\begin{proof}
    Let $\mathcal E:=\{|S_{t^*}|>\log n\}$.
    For $m\geq 0$, let $B_m$ be the event that, in the natural step-by-step coupling of the $Z$-chain and the $Y$-chain, the two chains use different updates at least once during the time interval from $t^*$ to $t^*+m$.

    We first record the consequence of the coupling on the good event
    $$
    \mathcal E^c\cap B^c_{\frac{n\log n}{k}}\cap
    \left\{
    \max\{\tau_1,\tau_2\}-t^*\leq \frac{n\log n}{k}
    \right\}.
    $$
    On this event, the number of initially unmarked cards is at most $\log n$,
    and the two chains use the same update at every step until both stopping
    times have occurred. Since marking schemes A and B mark exactly the same
    card whenever the same admissible update is used, the marked sets in the
    two chains evolve identically throughout this time interval. Therefore the
    two terminal marking times are equal, namely
    $$
    \tau_1=\tau_2,
    $$
    and the corresponding permutations also agree:
    $$
    Z_{\tau_1}=Y_{\tau_2}.
    $$
    Hence
    $$
    \{\tau_1\neq \tau_2\}\cup\{Z_{\tau_1}\neq Y_{\tau_2}\}
    \subset
    \mathcal E\cup
    \left(B_{\frac{n\log n}{k}}\cap \mathcal E^c\right)
    \cup
    \left\{
    \max\{\tau_1,\tau_2\}-t^*>\frac{n\log n}{k}
    \right\}.
    $$
    It follows that
    $$
    \mathbf P(\tau_1\neq \tau_2)
    \leq
    \mathbf P(\mathcal E)
    +
    \mathbf P\left(B_{\frac{n\log n}{k}}\cap \mathcal E^c\right)
    +
    \mathbf P\left(
    \max\{\tau_1,\tau_2\}-t^*>\frac{n\log n}{k}
    \right).
    $$
    The same bound also controls the coupling error between the stopped
    variables. By the coupling inequality,
    $$
    d_{\mathrm{TV}}\bigl(\mathcal L(Z_{\tau_1}),\mathcal L(Y_{\tau_2})\bigr)
    \leq
    \mathbf P(Z_{\tau_1}\neq Y_{\tau_2})
    $$
    and hence
    $$
    d_{\mathrm{TV}}\bigl(\mathcal L(Z_{\tau_1}),\mathcal L(Y_{\tau_2})\bigr)
    \leq
    \mathbf P(\mathcal E)
    +
    \mathbf P\left(B_{\frac{n\log n}{k}}\cap \mathcal E^c\right)
    +
    \mathbf P\left(
    \max\{\tau_1,\tau_2\}-t^*>\frac{n\log n}{k}
    \right).
    $$

    We now bound the three terms on the right-hand side. First, by the estimate
    on the number of cards not yet touched at time $t^*$, we have
    \begin{equation}\label{eq: St small}
        \mathbf P(\mathcal E)=\mathbf P(|S_{t^*}|>\log n)=n^{-\omega(1)}.
    \end{equation}
    Second, on the event $\mathcal E^c$, the number of unmarked cards is at most
    $\log n$ throughout the coupling. For simpler notation, we let \(s = |S_t|\) and \(|M_t| = n-s\). 
Then we compare the one-step transition rules of the two coupled chains and get
\begin{equation}
\begin{aligned}
d_{\mathrm{TV}}\!\left(P_{n,t},Q_{n,t}\right)= s C_t q_t 
= \frac{s}{n(n-1)}
   \frac{(n-s)k}{n-s+1}
   \prod_{i=2}^{k-1}
   \left(1-\frac{s-1}{n-i}\right) = \frac{sk}{n^2}\left(1+o(1)\right),
\end{aligned}
\tag{3.16}
\end{equation}
for \(k=o\!\left(n/(\log n)^4\right)\) and \(s=O(\log n)\).
    
We obtain by a union bound over the next
    $n\log n/k$ steps that
    $$
    \mathbf P\left(B_{\frac{n\log n}{k}}\cap \mathcal E^c\right)
    \leq
    \frac{n\log n}{k}\cdot
    \frac{k\log n}{n^2}(1+o(1))
    =
    n^{-1+o(1)}.
    $$
    Finally, we use the coupon-collector-type estimates for the two marking schemes. Fix a point in $S_{t^*}$. At each
subsequent step, this point is included in the chosen uniform $k$-cycle with probability $k/n$.
Hence, after the next $\left\lceil n\log n/k\right\rceil$ steps, the probability that this point
is still untouched is at most
\begin{equation}
    \left(1-\frac{k}{n}\right)^{\left\lceil n\log n/k\right\rceil}
    \leq \exp(-\log n)
    = n^{-1}.
\end{equation}
Taking a union bound over all points in $S_{t^*}$ gives
\begin{equation}\label{eq: coupon}
    \mathbf P\left(\tau-t^*>\frac{n\log n}{k},\ |S_{t^*}|\leq \log n\right)\leq\log n\cdot\left(1-\frac{k}{n}\right)^{\left\lceil n\log n/k\right\rceil}\leq\frac{\log n}{n}
    = n^{-1+o(1)}.
\end{equation}
As a result, we get
    $$
    \mathbf P\left(
    \tau_1-t^*>\frac{n\log n}{k},\,\mathcal E^c
    \right)
    \leq n^{-1+o(1)}\quad \text{and} \quad
    \mathbf P\left(
    \tau_2-t^*>\frac{n\log n}{k},\,\mathcal E^c
    \right)
    \leq n^{-1+o(1)}.
    $$
    Therefore
    $$
    \mathbf P\left(
    \max\{\tau_1,\tau_2\}-t^*>\frac{n\log n}{k}
    \right)
    \leq
    n^{-1+o(1)}+\mathbf P(\mathcal E)
    =
    n^{-1+o(1)}.
    $$

    Combining the preceding estimates yields
    $$
    \mathbf P(\tau_1\neq \tau_2)\leq n^{-1+o(1)}
    $$
    and
    $$
    d_{\mathrm{TV}}\bigl(\mathcal L(Z_{\tau_1}),\mathcal L(Y_{\tau_2})\bigr)
    \leq n^{-1+o(1)}.
    $$
    This completes the proof.
\end{proof}

%We then compare the two marking schemes by coupling the modified kernel $Q_{n,t}$ with the original $k$-cycle kernel $P_n$.

We next compare the terminal time of marking scheme B with the genuine hitting time $\tau$ of the original random $k$-cycle walk.

\begin{lemma}\label{lemma: Z,X gap}
    With notations defined in \Cref{subsec: marking schemes}, we have 
    $$d_{\mathrm{TV}}(\mathcal{L}(Z_{\tau_1}),\mathcal{L}(X_{\tau}))\leq \exp\left(-(\log\log n)^{1/2+o(1)}\right).$$
\end{lemma}

\begin{proof}
Let $\mathcal B$ be the event that, before time
$t^*+ n\log n/k$, some chosen $k$-cycle contains at least two elements of $S_{t^*}$ conditioned on $|S_{t^*}|\leq \log n$. Then, by a union bound,
$$
\mathbf P(\tau_1\neq \tau)\leq\mathbf P(\tau\leq t^*)+\mathbf P(|S_{t^*}|\geq \log n) +\mathbf P\left(\tau-t^*>\frac{n\log n}{k},\ |S_{t^*}|\leq \log n\right)+\mathbf P(\mathcal B).
$$
By \Cref{lemma::Self-reduction lemma}, we have
$$
    \mathbf P(\tau\leq t^*)
    \leq \mathbf P(\mathcal F^{t^*}(\varnothing)) =\left(1+(\log n)^{-1/2+o(1)}\right)
    \exp(-\gamma_{t^*}).
$$
Moreover,
\begin{align*}
    \gamma_{t^*}
    &=\exp\left(
    -\frac{\left(t^*-\frac{n\log n}{k}\right)k}{n}
    \right) \\
    &=\exp\left(
    \frac{1}{2}\log\log\log n
    -4\log\log\log\log n
    +o(1)
    \right) \\
    &=(\log\log n)^{1/2+o(1)}.
\end{align*}
Therefore,
\begin{equation}\label{eq: finish early}
    \mathbf P(\tau\leq t^*)
    \leq \exp\left(-(\log\log n)^{1/2+o(1)}\right).
\end{equation}
Also, by the definition of $t^*$ and the coupon-collector estimate for the number of untouched
points,
\begin{equation}
    \mathbf P(|S_{t^*}|\geq \log n)\leq n^{-\omega(1)}.
\end{equation}
We now condition on the event $|S_{t^*}|\leq \log n$.
By \eqref{eq: coupon}, we see that
$$\mathbf P\left(\tau-t^*>\frac{n\log n}{k},\ |S_{t^*}|\leq \log n\right)\leq\frac{\log n}{n}
    = n^{-1+o(1)}.$$

It remains to control the event that the touching scheme and the marking scheme differ before
time $t^*+ n\log n/k$. Conditional on $|S_{t^*}|\leq \log n$, at a single
step this can happen only if the chosen uniform $k$-cycle contains at least two elements of
$S_{t^*}$. The probability of this event is at most
\begin{equation}
    \binom{k}{2}\left(\frac{|S_{t^*}|}{n}\right)^2
    \leq \binom{k}{2}\frac{(\log n)^2}{n^2}.
\end{equation}
Summing over the next $\left\lceil n\log n/k\right\rceil$ steps, we obtain
\begin{equation}\label{eq: prob of B}
\mathbf P(\mathcal B)\leq\left\lceil\frac{n\log n}{k}\right\rceil\binom{k}{2}\frac{(\log n)^2}{n^2} \notag\leq \frac{k(\log n)^3}{n}\leq \frac{1}{\log n}.
\end{equation}
Combining the preceding estimates yields
\begin{equation}
    \mathbf P(\tau_1\neq \tau)
    \leq \exp\left(-(\log\log n)^{1/2+o(1)}\right).
\end{equation}

We first compare the initial laws of \(Z_{t^*}\) and \(X_{t^*}\). Recall that \(S_{t^*}\) denotes the set of cards which have not been touched by time \(t^*\). Conditional on the event \(S_{t^*}=M\), the random variable \(Z_{t^*}\) is sampled uniformly from \(\Omega_{t^*}(M)\), and hence has law \(\mu_{t^*}^M\). On the other hand, conditional on \(S_{t^*}=M\), the law of \(X_{t^*}\) is exactly the law of \(X_{t^*}\) conditioned on \(F(M)\). Therefore, by conditioning on \(S_{t^*}\), we obtain
$$
\begin{aligned}
    d_{\mathrm{TV}}\bigl(\mathcal L(Z_{t^*}),\mathcal L(X_{t^*})\bigr)
    &\leq
    \sum_{M\subseteq[n]}
    \mathbf P(S_{t^*}=M)\,
    d_{\mathrm{TV}}\bigl(\mu_{t^*}^M,\mathcal L(X_{t^*}\mid F(M))\bigr) \\
    &\leq
    (\log n)^{-1/2+o(1)}
    +
    \mathbf P(|S_{t^*}|>\log n).
\end{aligned}
$$
Here, in the second line we used \Cref{prop:key} for all \(M\) with \(|M|\leq \log n\), and the trivial bound \(d_{\mathrm{TV}}\leq 1\) for the remaining values of \(M\). By the coupon-collector estimate at time \(t^*\), we have
$$
\mathbf P(|S_{t^*}|>\log n)=n^{-\omega(1)}.
$$
Thus
$$
d_{\mathrm{TV}}\bigl(\mathcal L(Z_{t^*}),\mathcal L(X_{t^*})\bigr)
\leq
(\log n)^{-1/2+o(1)}+n^{-\omega(1)}.
$$

We now compare the laws at the stopping times. Consider a coupling of $(Z_t,X_t)$ such that $d_{TV}(\mathcal{L}(Z_{t^*}),\mathcal{L}(X_{t^*})) = \mathbf{P}(Z_{t^*} \neq X_{t^*})$ and $Z_t$, $X_t$ evolve according to the same transition kernel after time \(t^*\). The total variation distance between their laws cannot increase under this common Markov evolution. Then
% $$
% \begin{aligned}
%     d_{\mathrm{TV}}\bigl(\mathcal L(Z_{\tau_1}),\mathcal L(X_\tau)\bigr)
%     &\leq
%     d_{\mathrm{TV}}\bigl(\mathcal L(Z_\tau),\mathcal L(X_\tau)\bigr)
%     +\mathbf P(\tau_1\neq \tau) \\
%     &\leq
%     d_{\mathrm{TV}}\bigl(\mathcal L(Z_{t^*}),\mathcal L(X_{t^*})\bigr)
%     +\mathbf P(\tau_1\neq \tau) \\
%     &\leq
%     (\log n)^{-1/2+o(1)}
%     +n^{-\omega(1)}
%     +\exp\left(-(\log\log n)^{1/2+o(1)}\right) \\
%     &\leq
%     \exp\left(-(\log\log n)^{1/2+o(1)}\right).
% \end{aligned}
% $$

$$
\begin{aligned}
    &d_{\mathrm{TV}}\bigl(\mathcal L(Z_{\tau_1}),\mathcal L(X_\tau)\bigr)
    \leq
    d_{\mathrm{TV}}\bigl(\mathcal L(Z_{t^*}),\mathcal L(X_{t^*})\bigr)
    +\mathbf P(\tau_1\neq \tau)+  \mathbf P(\max(\tau_1,\tau) - t^*> n\log n/k) \\
    &\leq
    (\log n)^{-1/2+o(1)}
    +n^{-\omega(1)}
    +\exp\left(-(\log\log n)^{1/2+o(1)}\right)  + n^{-1+o(1)}\\
    &\leq
    \exp\left(-(\log\log n)^{1/2+o(1)}\right).
\end{aligned}
$$
\end{proof}

The next lemma shows that the auxiliary uniform process $\rho_t$ and the $Y$-chain have the same terminal marking time, and remain close at that time.

\begin{lemma}\label{lem: Y and rho}
    Under the natural coupling in which $(Y_t)_{t\ge t^*}$ and $(\rho_t)_{t\ge t^*}$ both use the same sampled move generated from $Q_{n,t}$ at each time and $Y_{t^*} = \rho_{t^*}$, one has $\tau_2=\tau_3$. Moreover,
    \[
    d_{\mathrm{TV}}(\mathcal L(Y_{\tau_2}), \mathcal L(\rho_{\tau_3})) \leq \frac{1+o(1)}{\log n}.
    \]
\end{lemma}

\begin{proof}
    Because we have coupled the two chains using the same sampled move $\sigma_t \sim Q_{n,t}$ at each time $t \geq t^*$ and set $Y_{t^*} = \rho_{t^*}$. The marked sets for these two chains are identical at any time and therefore $\tau_2 = \tau_3$.   
    Let $T = \frac{n \log n}{k} + t^*$. We know that
    \[
    \mathbf{P}\left(\tau_2  \geq T \right) \leq n^{-1+o(1)},\quad \mathbf{P}(|S_{t^*}| \geq \log n) \leq n ^{- \omega(1)}.
    \]
    Let $\mathcal{B}_T$ be the event that there exists some $1 \leq t \leq T - t^*$ such that $\sigma_{t+t^*}$ contains at least two elements from $[n] \setminus M_{t+t^*}$. Recall that $|[n] \setminus M_{t+t^*}| \leq |[n] \setminus M_{t^*}| = |S_{t^*}|$. By the same argument as in Lemma \ref{lemma: Z,X gap}, we get
    \[
    \mathbf{P} \left( \mathcal{B}_T \bigg| |S_{t^*}| \leq \log n \right) \leq \frac{1}{\log n}.
    \]
Recall that the processes \(Y_t\) and \(\rho_t\) are constructed on the same probability space, using the same \(Q_{n,t}\)-moves. Let \(\mathcal{B}_t\) be the bad event that the coupling between \(Y\) and \(\rho\) has failed by time \(t\). Then, on the event \(\mathcal{B}_{\tau_2}^c\), the two processes agree up to the marking time, and in particular \(Y_{\tau_2}=\rho_{\tau_3}\). Hence
$$
d_{\mathrm{TV}}\bigl(\mathcal L(Y_{\tau_2}),\mathcal L(\rho_{\tau_3})\bigr)
\leq
\mathbf P(Y_{\tau_2}\neq \rho_{\tau_3})
\leq
\mathbf P(\mathcal{B}_{\tau_2}).
$$
It remains to bound the probability of this bad event. We split according to whether the marking time \(\tau_2\) occurs before the deterministic time \(T\), and according to whether the number of initially unmarked cards is at most \(\log n\). Since \(\mathcal{B}_{\tau_2}\subseteq \{\tau_2\geq T\}\cup \mathcal{B}_T\), we have
$$
\begin{aligned}
    \mathbf P(\mathcal{B}_{\tau_2})
    &\leq
    \mathbf P(\tau_2\geq T)+\mathbf P(\mathcal{B}_T) \\
    &\leq
    \mathbf P(\tau_2\geq T)
    +\mathbf P\bigl(\mathcal{B}_T\,\big|\, |S_{t^*}|\leq \log n\bigr)
    +\mathbf P(|S_{t^*}|>\log n).
\end{aligned}
$$
By the preceding estimates, the three terms on the right-hand side satisfy
$$
\mathbf P(\tau_2\geq T)\leq n^{-1+o(1)}, \quad\mathbf P\bigl(\mathcal{B}_T\,\big|\, |S_{t^*}|\leq \log n\bigr)\leq \frac{1}{\log n},
$$
and
$$
\mathbf P(|S_{t^*}|\geq \log n)\leq n^{-\omega(1)}.
$$
Combining these estimates gives
$$
d_{\mathrm{TV}}\bigl(\mathcal L(Y_{\tau_2}),\mathcal L(\rho_{\tau_3})\bigr)\leq\frac{1}{\log n}+n^{-1+o(1)}+n^{-\omega(1)}\leq\frac{1+o(1)}{\log n}.
$$
This completes the proof.
\end{proof}

It remains to show that, when $k$ is even, the parity of the genuine hitting time $\tau$ is asymptotically balanced.

\begin{lemma}\label{lem:parity of the hitting time}
Suppose that \(k\) is even and \(2\leq k=o(n/(\log n)^4)\). Then 
$$ 
\mathbf P(\tau \text{ is even})=\frac12+O\left(\epsilon_n\right),
\qquad
\mathbf P(\tau \text{ is odd})=\frac12+O\left(\epsilon_n\right),
$$
where $\epsilon_n = \exp\left(-(\log \log n)^{1/2+o(1)}\right).$
In particular,
$$
d_{\mathrm{TV}}\left(
\mathbf P(\tau \text{ is even})U_{\mathfrak A_n}+\mathbf P(\tau \text{ is odd})U_{\mathfrak A_n^c},
U_{\mathfrak S_n}\right)=\exp\left(-(\log \log n)^{1/2+o(1)}\right).
$$
\end{lemma}

\begin{proof}
We first separate the exceptional event $\{\tau \leq t^*\}$. By \eqref{eq: finish early}, we already have $\mathbf{P}(\tau \leq t^*) \leq \epsilon_n.$ Let \(S_t\) be the set of cards not touched by time \(t\). By \eqref{eq: St small}, we  have that $\mathbf P(|S_{t^*}|>\log n)=n^{-\omega(1)}$.
Using the same argument as in Lemma \ref{lemma: Z,X gap}, we get that $$\mathbf{P}\left(
\mathcal B \bigg| |S_{t^*}|\leq \log n
\right)
\leq
O\left(\frac{k(\log n)^3}{n}\right)
=
o\left(\frac1{\log n}\right),$$
where $\mathcal B$ is the event that  at some step after $t^*$, the chosen $k$-cycle touches at least two elements which are still untouched at that time. So after time \(t^*\), with high probability the number of untouched cards decreases by one whenever it decreases. 
Therefore, with probability \(1-O(\epsilon_n)\), the process reaches a time
$$
\sigma:=\inf\{t\geq t^*: |S_t|=1\}
$$
before the hitting time \(\tau\).

On the event that \(\sigma<\tau\), let \(i\) be the unique card not touched by time \(\sigma\). After time \(\sigma\), the hitting time \(\tau\) is obtained by waiting until a sampled \(k\)-cycle contains \(i\). At each subsequent step this happens with probability \(k/n\), independently of the past. Hence
$\tau-\sigma$ has the geometric distribution on \(\{1,2,\ldots\}\) with success probability \(k/n\). Therefore,
$$
\left|\mathbf P(\tau-\sigma\text{ is even})-\mathbf P(\tau-\sigma\text{ is odd})\right|
=
\frac{k/n}{2-k/n}
=
O\left(\frac{k}{n}\right).
$$
Consequently, conditional on the history up to time \(\sigma\), the parity of \(\tau\) has bias at most \(O(k/n)\). Combining this with the exceptional probabilities above gives
\begin{equation}
    \begin{aligned}
        \left|
\mathbf P(\tau \text{ is even})-\frac12
\right|
&\leq
O\left(\frac{k}{n}\right)
+
\mathbf P(\tau\leq t^*)
+
\mathbf P(|S_{t^*}|>\log n)
+
\mathbf P\left(
\mathcal B\cap\{|S_{t^*}|\leq\log n\}
\right)=
O\left(\epsilon_n\right).
    \end{aligned}
\end{equation}

The same bound for \(\mathbf P(\tau\text{ is odd})\) follows since the two probabilities sum to one.

Finally, since $U_{\mathfrak S_n}=\frac12 U_{\mathfrak A_n}
+\frac12 U_{\mathfrak A_n^c}$, we have
\begin{equation}
d_{\mathrm{TV}}\left(
\mathbf P(\tau \text{ is even})U_{\mathfrak A_n}+\mathbf P(\tau \text{ is odd})U_{\mathfrak A_n^c},U_{\mathfrak S_n}\right)
=\left|\mathbf P(\tau \text{ is even})-\frac12\right|=\exp\left(-(\log \log n)^{1/2+o(1)}\right). 
\end{equation}
This completes the proof.
\end{proof}





\begin{proof}[Proof of \Cref{thm: hitting time mixing_main thm}]
We combine the uniformity of the auxiliary process \(\rho_t\) with the comparison estimates between the three auxiliary chains and the original chain.

First consider the case where \(k\) is odd. By the uniformity statement proved above, conditional on the marked set \(M_t\), the random permutation \(\rho_t\) is uniform on \(\Omega_t(S_t)\). At the stopping time \(\tau_3\), all elements have been marked, so \(S_{\tau_3}=\emptyset\). Since \(k\) is odd, the walk stays in \(\mathfrak A_n\), and therefore
$$
\rho_{\tau_3}\sim U_{\mathfrak A_n}.
$$
Hence, by the triangle inequality,
$$
d_{\mathrm{TV}}\bigl(\mathcal L(X_\tau),U_{\mathfrak A_n}\bigr)\leq
d_{\mathrm{TV}}\bigl(\mathcal L(X_\tau),\mathcal L(Z_{\tau_1})\bigr)
+d_{\mathrm{TV}}\bigl(\mathcal L(Z_{\tau_1}),\mathcal L(Y_{\tau_2})\bigr) 
+d_{\mathrm{TV}}\bigl(\mathcal L(Y_{\tau_2}),\mathcal L(\rho_{\tau_3})\bigr).
$$
By the estimates proved in the preceding lemmas,
$$
d_{\mathrm{TV}}\bigl(\mathcal L(X_\tau),\mathcal L(Z_{\tau_1})\bigr)
\leq
\exp\left(-(\log\log n)^{1/2+o(1)}\right),\,\,
d_{\mathrm{TV}}\bigl(\mathcal L(Z_{\tau_1}),\mathcal L(Y_{\tau_2})\bigr)
\leq
n^{-1+o(1)},
$$
and
$$
d_{\mathrm{TV}}\bigl(\mathcal L(Y_{\tau_2}),\mathcal L(\rho_{\tau_3})\bigr)
\leq
\frac{1+o(1)}{\log n}.
$$
Since
$$
n^{-1+o(1)}+\frac{1+o(1)}{\log n}
\leq
\exp\left(-(\log\log n)^{1/2+o(1)}\right),
$$
we conclude that
$$
d_{\mathrm{TV}}\bigl(\mathcal L(X_\tau),U_{\mathfrak A_n}\bigr)
\leq
\exp\left(-(\log\log n)^{1/2+o(1)}\right).
$$

It remains to treat the case where \(k\) is even. In this case each \(k\)-cycle is an odd permutation, so the parity of the walk alternates with time. Again, the uniformity of \(\rho_t\) conditional on the marked set implies that, at the marking time \(\tau_3\),
$$
\rho_{\tau_3}
\sim
\mathbf P(\tau_3\text{ is even})U_{\mathfrak A_n}
+
\mathbf P(\tau_3\text{ is odd})U_{\mathfrak A_n^c}.
$$
Since \Cref{lem: Y and rho} gives \(\tau_2=\tau_3\) under the coupling, and the previous comparison between \(Z_{\tau_1}\) and \(X_\tau\) includes the event \(\tau_1\neq \tau\), the same triangle inequality gives
\begin{multline}
d_{\mathrm{TV}}\left(
\mathcal L(X_\tau),
\mathbf P(\tau\text{ is even})U_{\mathfrak A_n}+
\mathbf P(\tau\text{ is odd})U_{\mathfrak A_n^c}
\right) \\\leq d_{\mathrm{TV}}\bigl(\mathcal L(X_\tau),\mathcal L(Z_{\tau_1})\bigr)
+d_{\mathrm{TV}}\bigl(\mathcal L(Z_{\tau_1}),\mathcal L(Y_{\tau_2})\bigr)
+d_{\mathrm{TV}}\left(\mathcal L(Y_{\tau_2}),
\mathbf P(\tau\text{ is even})U_{\mathfrak A_n}+\mathbf P(\tau\text{ is odd})U_{\mathfrak A_n^c}
\right).
\end{multline}
The last term is bounded by
$$
d_{\mathrm{TV}}\bigl(\mathcal L(Y_{\tau_2}),\mathcal L(\rho_{\tau_3})\bigr)
+
d_{\mathrm{TV}}\left(
\mathcal L(\rho_{\tau_3}),
\mathbf P(\tau\text{ is even})U_{\mathfrak A_n}
+
\mathbf P(\tau\text{ is odd})U_{\mathfrak A_n^c}
\right).
$$
The first part is at most \((1+o(1))/\log n\) by \Cref{lem: Y and rho}. For the second part, note that \(\rho_{\tau_3}\) has the parity mixture determined by \(\tau_3\), while the target mixture is determined by \(\tau\). Therefore this contribution is bounded by $\mathbf P(\tau_3\neq \tau)$. Since $\tau_3=\tau_2$ under the coupling, we have
$$
\mathbf P(\tau_3\neq \tau)\leq\mathbf P(\tau_2\neq \tau_1)+\mathbf P(\tau_1\neq \tau).
$$
By the comparison estimate between $\tau_1$ and $\tau_2$ in \Cref{lem: Z and Y}, together
with the bound
$$
\mathbf P(\tau_1\neq \tau)\leq\exp\left(-(\log\log n)^{1/2+o(1)}\right),
$$
we obtain
$$
\mathbf P(\tau_3\neq \tau)\leq n^{-1+o(1)}+\exp\left(-(\log\log n)^{1/2+o(1)}\right)=\exp\left(-(\log\log n)^{1/2+o(1)}\right).
$$
Thus
$$
d_{\mathrm{TV}}\left(
\mathcal L(X_\tau),
\mathbf P(\tau\text{ is even})U_{\mathfrak A_n}
+
\mathbf P(\tau\text{ is odd})U_{\mathfrak A_n^c}
\right)
\leq
\exp\left(-(\log\log n)^{1/2+o(1)}\right).
$$
Finally, by \Cref{lem:parity of the hitting time},
$$
d_{\mathrm{TV}}\left(
\mathbf P(\tau\text{ is even})U_{\mathfrak A_n}+\mathbf P(\tau\text{ is odd})U_{\mathfrak A_n^c},
U_{\mathfrak S_n}
\right)=\exp\left(-(\log\log n)^{1/2+o(1)}\right).
$$
Therefore, another application of the triangle inequality gives
$$
d_{\mathrm{TV}}\bigl(\mathcal L(X_\tau),U_{\mathfrak S_n}\bigr)
\leq
\exp\left(-(\log\log n)^{1/2+o(1)}\right).
$$
This completes the proof.
\end{proof}



